\section{Orienting the separations}
\label{sec:global-orientations}

Suppose that Theorem~\ref{thm:bilight} fails, and choose a
counterexample $G$ minimising $(v(G),e(G))$ lexicographically. Thus $G$
is $4$-bilight, has at least three vertices and $\rho(G)\geq-2$, and
has no $\Kminus$ minor. Simplicity gives $v(G)\geq9$. Throughout the
next two sections, $G$ denotes this minimum counterexample.

For a fragment $A$, write $R_A$ for $G[A\cup N_G(A)]$ rooted at
$N_G(A)$. We call a five-rooted side \emph{heavy} if its rooted
density is at least two, and \emph{low} if its density is at most one.
For a separation $(A,B)$, write $R_{A,B}=G[B]$, rooted at $A\cap B$.

\begin{lemma}\label{gl:reducible}
The graph $G$ has no reducible fragment.
\end{lemma}

\begin{proof}
Replace an inclusion-maximal reducible fragment by a clique or dart
rooted at its boundary. By Lemma~\ref{lem:reduction}, the resulting
minor is $4$-bilight and has density at least $\rho(G)$. The
definition of reducibility leaves at least three vertices outside the
fragment. The replacement strictly decreases the order: a clique
removes every interior vertex, and a dart retains one where there were
at least two. This contradicts minimality.
\end{proof}

The next lemma adapts the separation argument
of~\cite[Section~5]{DNR2026}. We give the details to account for
the density threshold of two and the final $\Kminus$ model.

\begin{lemma}\label{gl:orientations}
The following assertions hold.
\begin{enumerate}
\item If one side of a five-separation is heavy, the opposite side
is $4$-light.
\item For every five-separation with boundary $X$, exactly one side
is heavy. If $m=10-e(G[X])$, the heavy side has density at least $m+7$.
\item There is no bifragment $(S,T)$ with both boundaries of order
at most five, both densities positive, and every member with boundary
of order five heavy.
\end{enumerate}
\end{lemma}

\begin{proof}
We first establish an observation about small separations. Suppose
$(C,D)$ has order at most four, $R_{C,D}$ is $4$-light and has
nonpositive density, and
\[
 |(C\setminus D)\cup N_G(C\setminus D)|\geq3.
\]
There cannot be a five-separation $(A,B)$ with $B\subseteq D$ and
$R_{A,B}$ both $4$-light and heavy. To see this, remove from $C\cap D$
the vertices with no neighbour in $C\setminus D$, obtaining $C'$.
The rooted graph $H=R_{C',D}$ remains $4$-light and has nonpositive
density, by the root-forgetting observation
\cite[Observation~5.2]{DNR2026}: if a root set of order at most four
is reduced, each forgotten root contributes at most four incident
edges outside a given old fragment, offset by its density cost of four.
The roots of $H$ lie in $A$, so $(A\cap D,B)$ is a root separation
of $H$ with right side $R_{A,B}$. The helper model supplied there by
Theorem~\ref{thm:helper} contains one of the rooted graphs in
Lemma~\ref{lem:extraction}. That lemma gives a reducible fragment
of $H$, which is also reducible in $G$: at least $|C'|\geq3$
vertices remain outside it. This contradicts Lemma~\ref{gl:reducible}.

For (1), suppose that the side opposite a heavy side $R_{A,B}$ is
not $4$-light. It contains a positive fragment $L\subseteq A\setminus B$
with at most four neighbours. Bilightness implies that $R_{A,B}$ is
$4$-light, since a positive small-boundary fragment inside $B\setminus A$
would form a dense bifragment with $L$. Set
\[
 (C,D)=(L\cup N_G(L),\,V(G)\setminus L).
\]
The side $R_{C,D}$ is $4$-light and has nonpositive density: otherwise
its whole interior, or a positive small-boundary fragment inside it,
would again form a dense bifragment with $L$. Also
$|L\cup N_G(L)|\geq6$, since $L$ has positive density. Now
$B\subseteq D$, contradicting the preceding observation.

Let $a\geq b$ be the two densities at a five-separation with boundary
$X$. Then
\begin{equation}\label{gl:density-sum}
 a+b=\rho(G)+20-e(G[X])\geq m+8.
\end{equation}
In particular, $a\geq\lceil(m+8)/2\rceil$. If $b\geq2$, (1) makes
both sides $4$-light. The rooted clique bound
(Corollary~\ref{thm:rooted-clique}) applies to the side of density $a$;
Theorem~\ref{thm:helper} applies to the other. Join corresponding root
branch sets at their common root. The five resulting branch sets
are pairwise adjacent; the two helper branch sets remain in the second
open side and retain at least ten of their eleven possible adjacencies.
This gives a $\Kminus$ model. Thus
$b\leq1$, and \eqref{gl:density-sum} gives $a\geq m+7$, proving (2).

For (3), first consider two heavy five-fragments $S,T$ forming a
bifragment. By (1), both $R_S$ and $R_T$ are $4$-light. We claim that
five disjoint paths join their closed sides. Otherwise
Menger's theorem gives a separation $(C,D)$ of order at most four
containing $S\cup N_G(S)$ in $C$ and $T\cup N_G(T)$ in $D$.
Here $|C\setminus D|\geq3$, since $|S|\geq2$.

We verify that $R_{C,D}$ is $4$-light and has nonpositive density.
Discard from its root set the vertices with no neighbour in
$D\setminus C$, leaving a separation $(C,D')$. No vertex of $S$
can lie in $C\cap D'$, since all neighbours of $S$ lie in $C$.
Consequently $R_{C,D'}$ is a root side of order at most four in the
side opposite $R_S$. By (1) it has nonpositive density, equal to
that of $R_{C,D}$. Applying the same argument to every root side of
$R_{C,D}$ proves its $4$-lightness. The initial observation, applied
to $R_T$, gives a contradiction. The five paths therefore exist.

Trim these paths so that they run between $N_G(S)$ and $N_G(T)$,
with interiors outside both closed sides. The ends are distinct on
each five-element boundary; a shared boundary vertex uses a trivial
path. Apply the helper theorem in $R_S$ and the rooted clique bound
in $R_T$, using (2). Join corresponding root branch sets along these
paths. Their interiors are disjoint; assigning each to its corresponding
branch set preserves disjointness and connectivity. Again they give a
$\Kminus$ model.
This excludes two heavy five-fragments. A heavy five-fragment and a
positive fragment of boundary at most four contradict (1). Two
positive fragments of boundary at most four contradict bilightness.
\end{proof}

\begin{lemma}\label{gl:exact-density}
We have $e(G)=4v(G)-2$.
\end{lemma}

\begin{proof}
If $\rho(G)\geq-1$, delete an edge $uv$ to obtain $J$. Minimality
forces a dense bifragment $(S,T)$ in $J$, with both boundaries of
order at most four. Unless this already violates bilightness in $G$,
we may label the ends so that $u\in S$ and
$v\notin S\cup N_J(S)$. If $v\in T$, choose the labels also so that
$|N_G(S)|\geq|N_G(T)|$. Then $|N_G(S)|\leq5$ and
$\rho_G(S)=\rho_J(S)+1\geq2$, so $|S|\geq2$.

Put $S'=S\setminus\{u\}$. We have
\[
 N_G(S')\subseteq(N_G(S)\setminus\{v\})\cup\{u\},\qquad
 \rho_G(S')=\rho_G(S)+4-e_G(\{u\},N_G(S)).
\]
The pair $(S',T)$ is nonadjacent in $G$. If $|N_G(S)|\leq4$, both
members have boundary at most four and positive density, a
contradiction. Otherwise Lemma~\ref{gl:orientations}(2) gives
$\rho_G(S)\geq7$, hence $\rho_G(S')\geq6$. The density of $T$ is
positive, and at least two if its boundary grows to order five.
Lemma~\ref{gl:orientations}(3) again gives a contradiction.
\end{proof}

\section{Degree bounds and five-vertex cuts}
\label{sec:global-cuts}

\begin{lemma}\label{gl:five-connected}
The graph $G$ has minimum degree at least five and is five-connected.
\end{lemma}

\begin{proof}
If $d_G(v)\leq4$, the minor $J=G-v$ has density at least $-2$
and at least eight vertices. Any dense bifragment in $J$ with
boundaries of order at most four remains a bifragment in $G$. Its
members remain positive; if a boundary grows to order five, the
restored vertex adds at least one incident edge, so the density is
at least two. Lemma~\ref{gl:orientations}(3) excludes such a pair.
Thus $J$ is $4$-bilight, contrary to minimality.

Suppose that $G$ has connectivity $k\leq4$, and let $S$ be a minimum
cut, possibly empty. Each component of $G-S$ has neighbourhood $S$.
By bilightness some component $A$ has nonpositive density. The side
$R_A$ is internally $k$-connected, and its nonroots retain degree at
least five. For $k\leq2$, its connected interior gives a rooted
clique. For $k=3$, apply Lemma~\ref{lem:triangle}; for $k=4$,
apply Lemma~\ref{lem:dart}. In the last case $|A|\geq2$, since the
minimum degree is five. Another component together with $S$ has at
least six vertices, again by the degree bound. Hence $A$ is reducible,
contrary to Lemma~\ref{gl:reducible}.
\end{proof}

Write $t(xy)=|N_G(x)\cap N_G(y)|$ for an edge $xy$. Contracting this edge
changes the global density by
\begin{equation}\label{gl:contraction-density}
 \rho(G/xy)=\rho(G)+3-t(xy).
\end{equation}
A quasi-five-connected graph is $4$-bilight: any member of a dense
bifragment has boundary four and at least two vertices, and the
other member supplies at least two vertices on the opposite side of
that cut. Thus a quasi-five-connected contraction with $t(xy)\leq3$
would contradict minimality.

We will repeatedly use the following observation. Contracting an edge
of a five-connected graph gives a four-connected graph, and every
four-cut contains the merged vertex. Consequently, in a dense
bifragment of the contraction with boundaries of order at most four,
the merged vertex belongs to both boundaries. Its open sides lift to
positive five-fragments of the original graph, with both contracted
ends in each boundary. Their densities do not decrease; they remain
disjoint and nonadjacent. Each has at least two vertices, because a
singleton of boundary four has density zero in the contraction.

\begin{lemma}\label{gl:degree-six}
The minimum degree of $G$ is at least six.
\end{lemma}

\begin{proof}
Suppose that $d_G(v)=5$. If $N_G(v)$ is a clique, choose a component
$C$ outside $N_G[v]$. Such a component exists because $v(G)\geq9$,
and five-connectivity gives $N_G(C)=N_G(v)$. The five singleton
neighbours, $\{v\}$ and $C$ form a $\Kminus$ model.

Otherwise choose nonadjacent $a,b\in N_G(v)$ and set $J=G-v+ab$.
This is a minor: contract $va$, then delete the new edges at $a$
other than $ab$. It has density $-2$, so a dense bifragment $(Y,Z)$
with boundaries of order at most four exists in $J$. For either
member, say $Y$, put $k=e_G(\{v\},Y)$ and let $c$ be one if the
added edge $ab$ has an end in $Y$, and zero otherwise. Then
\begin{equation}\label{gl:split-density}
 \rho_G(Y)=\rho_J(Y)+k-c.
\end{equation}
Since $c=1$ implies $k\geq1$, both sets remain positive in $G$;
they are still disjoint and nonadjacent, and their boundaries have
order at most five. By Lemma~\ref{gl:orientations}(3), we may choose
$Y$ with boundary five and $\rho_G(Y)=1$. Its boundary
grows on restoring $v$, so $k\geq1$. Equation~\eqref{gl:split-density}
forces $\rho_J(Y)=k=c=1$. Exchange $a,b$ if necessary: $a$ is the
unique neighbour of $v$ in $Y$, and $b$ lies in the boundary of $Y$
in both graphs. Hence $|Y|\geq2$. The opposite side of $Y$ in $G$
contains $Z$, which has at least two vertices.

If $|Y|\geq3$, Lemma~\ref{lem:quasi} makes $G/va$
quasi-five-connected. As $ab$ is absent and $d_G(v)=5$, we have
$t(va)\leq3$, a contradiction. We are left with $Y=\{a,c\}$.
The equation
\[
 1=d_G(a)+d_G(c)-e(G[\{a,c\}])-8
\]
forces $ac\in E(G)$ and $d_G(a)=d_G(c)=5$. Their common neighbours
form a triple $R$, and their complete neighbourhoods are
\begin{equation}\label{gl:first-pair}
 N_G(a)=\{c,v\}\cup R,\qquad
 N_G(c)=\{a,b\}\cup R,\qquad |R|=3.
\end{equation}
In particular $vb$ is an edge.

Contract $ac$, which has exactly three common neighbours. By
minimality and the preceding contraction observation, a dense
bifragment lifts to two positive five-fragments. Choose a member
$T$ of density one, using Lemma~\ref{gl:orientations}(3). Lifting
increases its density by $|T\cap R|$, so $T$ avoids $R$. Both $a,c$
are in its boundary; thus $v,b\in T$, and $v$ is the unique
neighbour of $a$ in $T$. Its opposite side has at least two vertices.
If $|T|\geq3$, Lemma~\ref{lem:quasi} applies to $av$; the missing
edge $cv$ gives $t(av)\leq3$, again a contradiction. Therefore
$T=\{v,b\}$, and the same degree calculation gives
\begin{equation}\label{gl:second-pair}
 N_G(v)=\{a,b\}\cup D,\qquad
 N_G(b)=\{c,v\}\cup D,\qquad |D|=3.
\end{equation}
The four degree-five vertices induce the cycle $a,c,b,v,a$, and
$R,D$ avoid the cycle.

If $|R\cap D|\geq2$, the cycle has at most four external neighbours.
There is a vertex outside the cycle and its boundary, because
$v(G)\geq9$, contradicting five-connectivity. If $|R\cap D|=1$,
put $S=R\cup D$ and $U=\{a,c,b,v\}$. Here $|S|=5$ and
$\rho_G(U)=0$. Its opposite side $W$ has
\[
 \rho_G(W)=18-e(G[S])\geq8.
\]
In particular $W$ is nonempty. Its closed side rooted at $S$ is
internally five-connected, so the rooted clique theorem applies.
The connected branch sets $\{a,v\}$ and $\{c,b\}$ are adjacent to each
other and to every vertex of $S$. Together with the rooted clique
they give a $K_7$ model.

Finally suppose that $R\cap D=\varnothing$. Every five-cut $L$
containing $a,v$ must also contain $c,b$. Indeed, if $c\notin L$,
all of $R$ lies in its component of $G-L$ or in $L$. Then $a$ has
no neighbour in any other component, and removing $a$ from $L$
leaves a cut of order four. The argument for $b$ is the same, using
the neighbours of $v$.

Now $t(av)=0$. The contraction $G/av$ satisfies the density bound, so
it has a dense bifragment. Both lifted five-boundaries contain all
four cycle vertices, by the preceding paragraph. Each open side
meets both $R$ and $D$, because its boundary contains $a,v$.
One of these disjoint sides, say $B$, contains exactly one vertex
$r$ of $R$. It cannot be a singleton, since its image has positive
density and boundary four. If $|B|=2$, then $B=\{r,d\}$ for some
$d\in D$, and its boundary consists of the four cycle vertices
and one further vertex $z$. The vertex $r$ is adjacent to neither
$v$ nor $b$, so its neighbours would all lie in $\{a,c,d,z\}$,
contrary to the minimum degree five.

Thus $|B|\geq3$, its opposite side has at least two vertices,
and $a$ has the unique neighbour $r$ in $B$. Lemma~\ref{lem:quasi}
makes $G/ar$ quasi-five-connected. Also $t(ar)\leq3$, since $a$
has degree five and its neighbour $v$ is not adjacent to $r$.
This final contraction contradicts minimality. Every contraction
above leaves at least eight vertices, preserving the order condition.
\end{proof}

\begin{lemma}\label{gl:low-edge}
Every proper low five-fragment $A$ contains an endpoint of an edge
$e$ with $t(e)\leq3$.
\end{lemma}

\begin{proof}
Put $S=N_G(A)$ and $a=|A|$. The minimum degree six gives
\[
 \rho_G(A)=\sum_{v\in A}d_G(v)-e(G[A])-4a
 \geq2a-\binom a2,
\]
so $a\geq5$. Also
\[
 \sum_{v\in A}\bigl(d_G(v)+|N_G(v)\cap S|\bigr)
 =8a+2\rho_G(A)\leq8a+2.
\]
Choose $v\in A$ whose summand is at most eight. Suppose that every
edge at $v$ has at least four common neighbours. Let $H$ be
$G[N_G(v)]$ with the edges within $X_0=N_G(v)\cap S$ deleted.
Vertices of $H-X_0$ have degree at least four, and vertices in
$X_0$ have degree at least $5-|X_0|$. Moreover
$v(H)+|X_0|\leq8$.

The side $G[A\cup S]$, rooted at $S$, is internally five-connected.
Choose five disjoint paths from $S$ to $N_G(v)$, stopped when they
first reach $N_G(v)$. They avoid $v$, and the paths at vertices
of $X_0$ are trivial.
Their five ends form a set containing $X_0$. Apply
Lemma~\ref{lem:star} to $H$ rooted at these ends. Its central branch
set avoids $S$. Use this branch set and $\{v\}$ as helpers. Four
linkage paths become root branch sets. For the fifth, take the path
ending in the central branch set and delete its last vertex. The
trimmed path is nonempty, since its original end lies outside $X_0$
and hence outside $S$.
Both helpers are adjacent to the other four root branch sets and to
each other, and the trimmed branch set is adjacent to the central
helper. Thus at least ten of the eleven helper adjacencies are present.
Every branch set lies in $A\cup S$, and no artificial boundary edge
has been used.

The opposite side is internally five-connected and has density
at least $m+7$, where $m=10-e(G[S])$, by
Lemma~\ref{gl:orientations}(2). It supplies a rooted clique on $S$.
Joining corresponding root branch sets gives a $\Kminus$ model,
a contradiction.
\end{proof}

The following elementary argument will exclude the remaining
five-cuts. A proper five-fragment is called an $F$-fragment, for a
family $F$ of edges, if its boundary contains both ends of an edge
in $F$. An $F$-end is an $F$-fragment containing no $F$-fragment as
a strict subset.

\begin{lemma}\label{gl:end}
Let $J$ be five-connected, with every proper five-fragment of order
at least four. If $A$ is an $F$-end, no edge $e\in F$ can meet $A$,
lie in $A\cup N_J(A)$ and have both ends in a five-cut of $J$.
\end{lemma}

\begin{proof}
Put $S=N_J(A)$ and $C=V(J)\setminus(A\cup S)$. Suppose a five-cut
$T$ contains both ends of such an edge $e$, and partition $J-T$
into two nonempty unions of components $B,D$. Write
\[
 \begin{array}{lll}
 p=|A\cap T|,&s=|S\cap T|,&r=|C\cap T|,\\
 u=|S\cap B|,&w=|S\cap D|.&
 \end{array}
\]
Then $p+s+r=u+s+w=5$, $p\geq1$ and $p+s\geq2$.

If $A\cap B\ne\varnothing$, its neighbourhood lies in
$(A\cap T)\cup(S\cap T)\cup(S\cap B)$. This set has order at least
five, by connectivity. If it has order five, it equals the boundary
and contains both ends of $e$, giving an $F$-fragment strictly
inside $A$. Its opposite side contains
$C$, and strict containment follows from $p\geq1$. Therefore
\[
 p+s+u\geq6,\qquad p>w,\qquad u>r.
\]
The opposite corner $C\cap D$ has boundary contained in a set of
order $r+s+w=5+r-u<5$, so it is empty. Symmetrically, if
$A\cap D\ne\varnothing$, then $p>u$, $w>r$ and $C\cap B$ is empty.
In each application, a nonempty opposite corner ensures that
the separation is proper.

If both $A$-corners are nonempty, then $C\subseteq T$, so
$r=|C|\geq4$, contrary to $p+s\geq2$. If just $A\cap B$ is
nonempty, then $D=S\cap D$, giving $w=|D|\geq4$. Since $p>w$,
we have $p=5$ and $r=s=0$. Now $C=C\cap B$ has all its neighbours
in $S\cap B$, of order $u=5-w\leq1$, again impossible. The other
one-corner case is symmetric.

If neither $A$-corner is nonempty, then $A\subseteq T$ and
$p=|A|\geq4$. Since $|C|\geq4$ and $r\leq1$, a $C$-corner is
nonempty, say $C\cap B$. Its boundary has order at most $r+s+u$,
so connectivity gives $u\geq p\geq4$. The other $C$-corner must
be empty, or similarly $w\geq4$. But then $D=S\cap D$, and again
$w=|D|\geq4$. This contradicts $u+w\leq5$.
\end{proof}

\begin{lemma}\label{gl:six-connected}
The graph $G$ is six-connected.
\end{lemma}

\begin{proof}
It is already five-connected. We first show that every proper
five-fragment has at least four vertices. A singleton would have
degree five. For a fragment of order $a=2$ or $3$, the degree bound
and the possible incident edges give
\[
 2a-\binom a2\leq\rho_G(A)\leq\binom a2+a.
\]
Its density is therefore three if $a=2$, and between three and six
if $a=3$. Both contradict Lemma~\ref{gl:orientations}(2), which
allows only densities at most one or at least seven. Also, every
low five-fragment is proper: an empty opposite side would give
density $18-e(G[N_G(A)])\geq8$.

Let $F$ be the edges with at most three common neighbours. Every
edge in $F$ has both ends in the boundary of a five-fragment of
density one. Indeed, its contraction preserves the density bound
by \eqref{gl:contraction-density}, so is not $4$-bilight. Its dense
bifragment lifts to two positive five-fragments, as described above;
Lemma~\ref{gl:orientations}(3) forces one of their densities to equal one.

Suppose a five-cut exists. It has a low side, so
Lemma~\ref{gl:low-edge} supplies an edge of $F$. The preceding
paragraph supplies a low $F$-fragment. Choose an inclusion-minimal
one, $A$. In fact $A$ is an $F$-end without any density restriction.
For a smaller $F$-fragment $B\subsetneq A$, a heavy $B$ would be
disjoint and nonadjacent to the heavy opposite side of $A$, violating
Lemma~\ref{gl:orientations}(3). Thus $B$ would be low, contradicting
the choice of $A$.

Lemma~\ref{gl:low-edge} now supplies an edge $e\in F$ meeting $A$.
Its other end lies in $A\cup N_G(A)$, and both ends lie in a
five-cut by the earlier contraction argument. This contradicts
Lemma~\ref{gl:end} and excludes every five-cut.
\end{proof}

\section{The final degree cases}
\label{sec:global-terminal}

The degree-six argument adapts the final part of the proof of
Norin and Totschnig in~\cite[Section~3]{NT2025}; six-connectivity
excludes the separator outcome directly. It uses
the two-paths theorem of Robertson, Seymour and
Thomas~\cite[Theorem~(2.4)]{RST1993}, in the form also
quoted in \cite[Theorem~13]{NT2025}. For four distinct vertices in
a specified cyclic order, it gives disjoint paths joining opposite
pairs, a separation of order at most three with all terminals on
one side and at least two vertices in the other open side, or a
drawing in a disc with the terminals on its boundary in that order.

\begin{lemma}\label{gl:degree-six-terminal}
Let $H$ be six-connected, with at least nine vertices and no
$\Kminus$ minor. If every edge has at least four common neighbours
and $H$ has a vertex of degree six, then $e(H)\leq4v(H)-9$.
\end{lemma}

\begin{proof}
Let $d_H(v)=6$. The complement of $H[N_H(v)]$ is a matching, since
each neighbour of $v$ has at least four neighbours there. Label the
six neighbours $u_1,w_1,u_2,w_2,u_3,w_3$ so that the only possible
nonedges are $u_iw_i$.

There cannot be disjoint paths joining two different pairs $u_i,w_i$
with interiors outside $N_H[v]$: contracting the interior of each
path into one of its ends supplies two of the three possibly missing
adjacencies. Together with $\{v\}$, the six resulting branch sets give
$\Kminus$. In particular all three pairs are nonedges. Otherwise
an edge joining one pair and a path for another pair through an exterior
component would suffice. Such a component exists and is adjacent
to all six neighbours of $v$, by six-connectivity.

Some pair, say $u_1,w_1$, has no common neighbour outside $N_H[v]$.
Otherwise choose an exterior common neighbour $x_i$ for each pair.
Two distinct choices give the forbidden disjoint length-two paths,
so all choices coincide at a vertex $x$. There is a component $C$
outside $N_H[v]\cup\{x\}$, because $v(H)\geq9$. Its boundary is
contained in $N_H(v)\cup\{x\}$ and has at least six vertices.
Thus at least five neighbours of $v$ meet $C$, giving a path for
one pair through $C$ and a disjoint path for another through $x$.

Apply the two-paths theorem to
$H'=H-\{v,u_1,w_1\}$, with ordered terminals $u_2,u_3,w_2,w_3$.
The two paths are excluded above. A separation $(A,B)$ in the
second outcome lifts to
\[
 (A\cup\{v,u_1,w_1\},\ B\cup\{u_1,w_1\})
\]
in $H$, of order at most five. It is proper: $v$ lies in the first
open side, and the second contains $B\setminus A$. No neighbour
of $v$ lies in that second open side. This contradicts
six-connectivity.

Hence $H'$ has the disc drawing from the third outcome. Its four
distinct boundary terminals give $e(H')\leq3v(H')-7$.
Every vertex outside $N_H[v]$
meets the deleted triple at most once, by the choice of $u_1,w_1$.
Each of the other four neighbours of $v$ meets it three times,
and the triple induces two edges. With $n'=v(H')=v(H)-3$,
\[
 e(H)\leq(3n'-7)+(n'-4)+12+2=4v(H)-9.
\]
\end{proof}

\begin{proof}[Proof of Theorem~\ref{thm:bilight}]
The minimum counterexample $G$ is six-connected by
Lemma~\ref{gl:six-connected}. Every edge has at least four common
neighbours: otherwise contracting it gives a five-connected, hence
$4$-bilight, smaller graph of density at least $-2$. The average
degree is $8-4/v(G)<8$, whereas Lemma~\ref{gl:degree-six} gives
minimum degree at least six. Thus some vertex has degree six or seven.

Degree six contradicts Lemma~\ref{gl:degree-six-terminal}, since
$e(G)=4v(G)-2$. For a degree-seven vertex $v$, put $J=G[N_G(v)]$.
The common-neighbour bound gives $\delta(J)\geq4$. Choose a
component $C$ outside $N_G[v]$, which exists since $v(G)\geq9$.
Its boundary lies in $N_G(v)$ and has at least six vertices.
Contract $C$ to a vertex and delete the other exterior components.
The resulting nine-vertex minor consists of $J$, the vertex $v$
adjacent to all of $J$, and a nonadjacent second vertex adjacent to
at least six vertices of $J$. Lemma~\ref{lem:quotient} gives a
$\Kminus$ model in this minor. Replacing the contracted vertex by
the connected set $C$ lifts the model to $G$.

Both degree cases are impossible, so no counterexample exists.
\end{proof}

\begin{proof}[Proofs of Corollary~\ref{thm:five-connected} and
Theorem~\ref{thm:colour}]
Every five-connected graph is $4$-bilight, giving
Corollary~\ref{thm:five-connected}. For the colouring theorem, suppose
a $\Kminus$-minor-free graph is not six-colourable, and choose such
a graph $G$ minor-minimal. Then its chromatic number is at least
seven and every proper minor is six-colourable.
The critical-graph reduction~\cite[Theorem~1.6]{DNR2026} gives
seven-connectivity and $e(G)\geq4v(G)-2$. The graph is therefore
$4$-bilight, and Theorem~\ref{thm:bilight} supplies the forbidden
minor.
\end{proof}
